% eksperymentalne tłumaczenie na włoski

%

\subsection{Okręgi dopisane} % lang-pl
\index{okrąg dopisany}% % lang-pl
\subsection{Circonferenze exinscritte} % lang-it
\index{circonferenza exinscritta}% % lang-it

Oprócz okręgu wpisanego w trójkąt są jeszcze trzy inne okręgi styczne do wszystkich trzech prostych, na których leżą boki trójkąta. % lang-pl
Nazywamy je okręgami dopisanymi (\emph{excircle} albo \emph{escribed circle} po angielsku), pisze o nich Audin \cite[s. 98]{audin_2003}. % lang-pl
Oltre alla circonferenza inscritta in un triangolo esistono altre tre circonferenze tangenti a tutte e tre le rette contenenti i lati del triangolo. % lang-it
Sono dette circonferenze exinscritte (\emph{excircle} oppure \emph{escribed circle} in inglese); Audin ne parla a \cite[p. 98]{audin_2003}. % lang-it

Środek okręgu wpisanego znajduje się na przecięciu dwusiecznych wewnętrznych kątów trójkąta. % lang-pl
Analogicznie środki okręgów dopisanych znajdują się na przecięciu dwusiecznej kąta wewnętrznego przy ustalonym wierzchołku i dwusiecznych kątów zewnętrznych przy pozostałych wierzchołkach. % lang-pl
Il centro della circonferenza inscritta è il punto d'incontro delle bisettrici interne del triangolo. % lang-it
Analogamente, i centri delle circonferenze exinscritte si trovano all'intersezione della bisettrice dell'angolo interno in un vertice fissato e delle bisettrici degli angoli esterni negli altri due vertici. % lang-it
% TODO: https://upload.wikimedia.org/wikipedia/commons/thumb/6/6f/Incircle_and_Excircles.svg/960px-Incircle_and_Excircles.svg.png
Ich promienie można łatwo wyznaczyć przez porównanie pól trójkątów, co prowadzi do wzoru: % lang-pl
I loro raggi si determinano facilmente confrontando aree di triangoli, ottenendo la formula % lang-it
\begin{equation}
	r_a = \frac{S_{ABC}}{p - a},
\end{equation}
oraz dwóch kolejnych dla $r_b$ i $r_c$, jak pokazuje Zetel \cite[s. 18]{zetel_2020}. % lang-pl
Tutaj $p$ oznacza tak jak zawsze połowę obwodu. % lang-pl
Czynnik $p - a$ sugeruje związek ze wzorem Herona i tak istotnie jest, mamy % lang-pl
e le analoghe formule per $r_b$ e $r_c$, come mostra Zetel \cite[p. 18]{zetel_2020}. % lang-it
Qui $p$ indica, come di consueto, il semiperimetro. % lang-it
Il fattore $p-a$ suggerisce un legame con la formula di Erone, e infatti % lang-it
\begin{equation}
	S_{ABC}^2 = r \cdot r_a \cdot r_b \cdot r_c.
\end{equation}

Zależności między czterema promieniami jest bez liku, mamy też: % lang-pl
Le relazioni tra questi quattro raggi sono numerosissime; ad esempio % lang-it
\begin{equation}
	\frac{1}{r} = \frac {1} {r_a} + \frac {1} {r_b} + \frac {1} {r_c}.
\end{equation}

% TODO: Wzory na promienie okręgów wpisanych, dopisanych.
\todofoot{$4R = r_a + r_b + r_c - r$} % Coxeter, s. 13; dopisz też bend okręgi Kartezjusza